% GENERATED FILE. Edit canonical lessons, then run npm run generate:books.
% canonical-source-hash: fnv1a64:f1a397b8f3276007
% canonical-lessons: GE-C279-ereignis, GE-C279-ereignen, GE-C279-stattfinden, GE-C279-vorkommen, GE-C279-vorfall

\chapter{Was geschah}
\label{ch:ge-was-geschah}
\hlchaptermodality{279}

\begin{chapteropening}
\textbf{By the end of this chapter:} \emph{I can say that something took place: an event, to occur, to take place, to happen and an incident.}

Complete the last lesson of chapter 279: I can say that something took place: an event, to occur, to take place, to happen and an incident.
\end{chapteropening}

\section[das Ereignis]{das Ereignis — the event}

\label{lesson:GE-C279-ereignis}



\begin{quote}
Before the new one: say the German for to sign, then the German for the conference.
\end{quote}

\subsection*{You'll want to know: das Ereignis}
\textbf{das Ereignis} — \textquotedblleft{}the event\textquotedblright{}.

\begin{quote}
Die Hochzeit war das große Ereignis des Jahres.
\end{quote}

\begin{quote}
\emph{The wedding was the big event of the year.}
\end{quote}

Plural: \emph{die Ereignisse}.

\begin{cousinweb}
\emph{Ereignis} is usually traced to Middle High German \emph{eröugen} \textquotedblleft{}to bring before the eyes\textquotedblright{}, built on \emph{Auge} \textquotedblleft{}eye\textquotedblright{}; the spelling with \emph{-ei-} came later. EtymDB lists no parent.
\end{cousinweb}

\begin{grammarlens}[title={What you've built}]
One more word for telling, explaining and describing at B1.
\end{grammarlens}

\subsection*{Your turn}
\begin{itemize}
\raggedright
  \item \emph{Say it:} \emph{das Ereignis}
  \item \emph{Say it:} \emph{das Ereignis}, once more
  \item \emph{Say it:} a sentence of your own with \emph{das Ereignis}
  \item \emph{Recall:} say the German for to sign, then the German for the conference, then say \emph{das Ereignis} again
\end{itemize}

\begin{tcolorbox}[breakable,colback=teal!4,colframe=teal!35!black,title={Before you move on}]
What does das Ereignis mean? (\textquotedblleft{}The event\textquotedblright{}.) Say it once more.
\end{tcolorbox}
\section[ereignen]{ereignen — to occur; reflexive: sich ereignen}

\label{lesson:GE-C279-ereignen}



\begin{quote}
Before the new one: say the German for the conference, then the German for the event.
\end{quote}

\subsection*{You'll want to know: ereignen}
\textbf{ereignen} — \textquotedblleft{}to occur; reflexive: sich ereignen\textquotedblright{}.

\begin{quote}
Der Unfall ereignete sich am frühen Morgen.
\end{quote}

\begin{quote}
\emph{The accident occurred early in the morning.}
\end{quote}

Only reflexive, and only in the third person: \emph{es ereignete sich}.

\begin{cousinweb}
\emph{ereignen} is usually traced to Middle High German \emph{eröugen} \textquotedblleft{}to show, to bring before the eyes\textquotedblright{}, built on \emph{Auge} \textquotedblleft{}eye\textquotedblright{}; EtymDB lists no parent.
\end{cousinweb}

\begin{grammarlens}[title={What you've built}]
One more word for telling, explaining and describing at B1.
\end{grammarlens}

\subsection*{Your turn}
\begin{itemize}
\raggedright
  \item \emph{Say it:} \emph{ereignen}
  \item \emph{Say it:} \emph{ereignen}, once more
  \item \emph{Say it:} a sentence of your own with \emph{ereignen}
  \item \emph{Recall:} say the German for the conference, then the German for the event, then say \emph{ereignen} again
\end{itemize}

\begin{tcolorbox}[breakable,colback=teal!4,colframe=teal!35!black,title={Before you move on}]
What does ereignen mean? (\textquotedblleft{}To occur; reflexive: sich ereignen\textquotedblright{}.) Say it once more.
\end{tcolorbox}
\section[stattfinden]{stattfinden — to take place}

\label{lesson:GE-C279-stattfinden}



\begin{quote}
Before the new one: say the German for the event, then the German for to occur.
\end{quote}

\subsection*{You'll want to know: stattfinden}
\textbf{stattfinden} — \textquotedblleft{}to take place\textquotedblright{}.

\begin{quote}
Das Konzert fand im Park statt.
\end{quote}

\begin{quote}
\emph{The concert took place in the park.}
\end{quote}

Separable: \emph{es findet statt}, \emph{es hat stattgefunden}. For planned events.

\begin{cousinweb}
\emph{stattfinden} joins the old noun \emph{Statt} \textquotedblleft{}place\textquotedblright{} and \emph{finden} \textquotedblleft{}to find\textquotedblright{}, a cousin of English \emph{find} (EtymDB): to find its place.
\end{cousinweb}

\begin{grammarlens}[title={What you've built}]
One more word for telling, explaining and describing at B1.
\end{grammarlens}

\subsection*{Your turn}
\begin{itemize}
\raggedright
  \item \emph{Say it:} \emph{stattfinden}
  \item \emph{Say it:} \emph{stattfinden}, once more
  \item \emph{Say it:} a sentence of your own with \emph{stattfinden}
  \item \emph{Recall:} say the German for the event, then the German for to occur, then say \emph{stattfinden} again
\end{itemize}

\begin{tcolorbox}[breakable,colback=teal!4,colframe=teal!35!black,title={Before you move on}]
What does stattfinden mean? (\textquotedblleft{}To take place\textquotedblright{}.) Say it once more.
\end{tcolorbox}
\section[vorkommen]{vorkommen — to happen, to occur now and then; to seem}

\label{lesson:GE-C279-vorkommen}



\begin{quote}
Before the new one: say the German for to occur, then the German for to take place.
\end{quote}

\subsection*{You'll want to know: vorkommen}
\textbf{vorkommen} — \textquotedblleft{}to happen, to occur now and then; to seem\textquotedblright{}.

\begin{quote}
So etwas kommt manchmal vor.
\end{quote}

\begin{quote}
\emph{Something like that happens sometimes.}
\end{quote}

Separable, Perfekt with \emph{sein}: \emph{Das ist schon oft vorgekommen}. \emph{Das kommt mir bekannt vor} — that seems familiar to me.

\begin{cousinweb}
\emph{vorkommen} joins \emph{vor} \textquotedblleft{}forward\textquotedblright{} and \emph{kommen} \textquotedblleft{}to come\textquotedblright{}, from Old High German \emph{kweman}, Proto-Germanic \emph{kwemaną} (EtymDB): to come forward.
\end{cousinweb}

\begin{grammarlens}[title={What you've built}]
One more word for telling, explaining and describing at B1.
\end{grammarlens}

\subsection*{Your turn}
\begin{itemize}
\raggedright
  \item \emph{Say it:} \emph{vorkommen}
  \item \emph{Say it:} \emph{vorkommen}, once more
  \item \emph{Say it:} a sentence of your own with \emph{vorkommen}
  \item \emph{Recall:} say the German for to occur, then the German for to take place, then say \emph{vorkommen} again
\end{itemize}

\begin{tcolorbox}[breakable,colback=teal!4,colframe=teal!35!black,title={Before you move on}]
What does vorkommen mean? (\textquotedblleft{}To happen, to occur now and then; to seem\textquotedblright{}.) Say it once more.
\end{tcolorbox}
\section[der Vorfall]{der Vorfall — the incident}

\label{lesson:GE-C279-vorfall}



\begin{quote}
Before the new one: say the German for to take place, then the German for to happen.
\end{quote}

\subsection*{You'll want to know: der Vorfall}
\textbf{der Vorfall} — \textquotedblleft{}the incident\textquotedblright{}.

\begin{quote}
Nach dem Vorfall rief er die Polizei.
\end{quote}

\begin{quote}
\emph{After the incident he called the police.}
\end{quote}

Plural: \emph{die Vorfälle}.

\begin{cousinweb}
\emph{Vorfall} joins \emph{vor} and \emph{Fall} \textquotedblleft{}a fall, a case\textquotedblright{}, from Old High German \emph{val} (EtymDB): something that falls out before you.
\end{cousinweb}

\begin{grammarlens}[title={What you've built}]
One more word for telling, explaining and describing at B1.
\end{grammarlens}

\subsection*{Your turn}
\begin{itemize}
\raggedright
  \item \emph{Say it:} \emph{der Vorfall}
  \item \emph{Say it:} \emph{der Vorfall}, once more
  \item \emph{Say it:} a sentence of your own with \emph{der Vorfall}
  \item \emph{Recall:} say the German for to take place, then the German for to happen, then say \emph{der Vorfall} again
\end{itemize}

\begin{tcolorbox}[breakable,colback=teal!4,colframe=teal!35!black,title={Before you move on}]
What does der Vorfall mean? (\textquotedblleft{}The incident\textquotedblright{}.) Say it once more.
\end{tcolorbox}
